% Second-question draft. Geometry and parameters follow repository main.
\subsection{问题二模型的建立与求解}\label{sec:q2}

\subsubsection{首测信息与候选源区域}

问题一将多次示向度转化为定位区域，并给出了最小包围圆半径 $R^*$ 对位置不确定性的刻画。问题二进一步要求：在只有首测信息时，主动选择第二检测点，使再次测向后的位置不确定性尽量减小。由于首测只提供方向区间而不提供距离，不能把示向射线上的某一点当作真实源位置，也不能直接按未知源坐标构造直角交会。

沿用第一问记号，首个检测点为 $S_1$，示向度为 $\theta_1$，对应误差角域为 $W_1$。以 $\overline B(c,r)$ 表示闭圆盘，目标区域为 $\Omega=\overline B((0,0),1800)$。普通示向响应说明源与首站距离大于 $5$ m，并且不超过有效接收半径的上限 $1500$ m，因此首测精确可行集为
\begin{equation}\label{eq:q2-source}
\mathcal C_1=\bigl(\Omega\cap W_1\cap\overline B(S_1,1500)\bigr)\setminus\overline B(S_1,5).
\end{equation}
为使用第一问的凸多边形运算，计算时省略 $5$ m 内圆排除，并用外切正多边形包络圆边界，得到 $P_1\supseteq\mathcal C_1$。这里 $P_1$ 表示源可能位于何处，下面的 $K_0$ 则表示第二检测点可以选在何处，二者的几何含义不同。

\subsubsection{保证再次接收的第二点候选区域}

全向源的有效接收半径 $\rho\in[1000,1500]$ m 未知。为避免移动后完全失去信号，先要求第二点 $q$ 到每个候选源位置都不超过最小接收半径，定义
\begin{equation}\label{eq:q2-safe}
\begin{aligned}
K_0(P_1)&=\left\{q:\max_{X\in P_1}\|q-X\|_2\le1000\right\}\\
&=\bigcap_{v\in\operatorname{Vert}(P_1)}\overline B(v,1000).
\end{aligned}
\end{equation}
第二个等号不需要遍历整个区域。事实上，对任意 $X=\sum_i\alpha_i v_i\in P_1$，其中 $\alpha_i\ge0$ 且 $\sum_i\alpha_i=1$，有
\[
\|q-X\|_2\le\sum_i\alpha_i\|q-v_i\|_2\le\max_i\|q-v_i\|_2.
\]
因此只要逐顶点距离检查通过，就能覆盖所有候选位置。又因真实源位于 $P_1$ 内且 $\rho\ge1000$ m，任意 $q\in K_0$ 均保证再次接收到信号：距离大于 $5$ m 时返回示向度，不超过 $5$ m 时返回近场响应，可直接进行光学定位。

第一问的最小包围圆还给出候选域的非空判据：
\begin{equation}\label{eq:q2-nonempty}
K_0(P_1)\ne\varnothing\quad\Longleftrightarrow\quad R^*(P_1)\le1000.
\end{equation}
此时 $C^*(P_1)$ 就是一个可行第二点。为直观说明首测后存在可选空间，令 $u(\varphi_1)=(\cos\varphi_1,\sin\varphi_1)$，完整首测扇形可被下述圆包络：
\[
c_1=S_1+\frac{1500}{2\cos\delta}u(\varphi_1),\qquad R_1=\frac{1500}{2\cos\delta}.
\]
沿任一射线，到 $c_1$ 的距离平方关于径向距离为凸函数，检查扇形顶点和外弧端点即可验证包络。当 $\delta=\pi/180$ 时 $R_1\approx750.114$ m，故 $\overline B(c_1,1000-R_1)$ 给出一个明确的保证再接收区域；加入目标圆先验只会减小源候选集。实际离散计算仍以式\eqref{eq:q2-safe}逐顶点核验为准。

$K_0$ 是\textbf{保守的保证再接收域}。首测有信号还说明 $\rho\ge\|X-S_1\|$，充分利用这一联合信息可能允许更多第二点；本文选用易于验证的式\eqref{eq:q2-safe}，不将它称为最大可行域。实现中将各半径 $1000$ m 圆盘替换为内接多边形以生成候选；内近似为空不能单独据此判定真实圆弧交集为空，仍须使用式\eqref{eq:q2-nonempty}检验。

\subsubsection{交会几何与候选点评分}

再次接收只是可行性要求，还需避免两次视线接近平行。对假设源位置 $X$，记两站到源的距离为 $r_1=\|X-S_1\|$、$r_2=\|X-q\|$，两条视线的锐交会角为 $\psi$。对测向方程作局部线性化，横向位置误差的尺度约为 $r_i\tan\delta$，两条约束法向量所组成矩阵的行列式绝对值为 $|\sin\psi|$，由此得到误差代理尺度
\begin{equation}\label{eq:q2-error}
E(q,X)=\frac{(r_1+r_2)\tan\delta}{|\sin\psi(q,X)|}.
\end{equation}
该式解释了交会角较大通常有利、近平行方向误差放大的原因，但它是一阶近似，并不是原非线性定位区域半径的严格上界。未知 $X$ 时，应对多个相容位置评价，而非只对首测射线上的一个估计点设计交会。

根据首测可得的 $P_1$ 和 $K_0$ 构造有限候选检测点集 $Q$，同时从 $P_1$ 的顶点及最小圆心构造有限假设集 $H$。采用定位误差代理与移动时间的组合评分
\begin{equation}\label{eq:q2-score}
J(q)=\max_{X\in H}\frac{(\|X-S_1\|+\|X-q\|)\tan\widetilde\delta}
{\max\{|\sin\psi(q,X)|,0.02\}}+\lambda\frac{\|q-S_1\|}{5},
\end{equation}
其中速度为 $5$ m/s，$\lambda=0.02$ m/s，使两项均具有长度量纲。$\widetilde\delta=1.005\pi/180$ 是实现中使用的保守半角；相对于题面 $1^\circ$，额外的 $0.005^\circ$ 用于容纳接口两位小数示向的舍入，构造 $P_1$ 时使用同一半角。分母下限 $0.02$ 用于控制近平行情况的数值发散。这些是评分与实现参数，不是新的物理假设。

如果假设源距第二点不超过 $5$ m，评分使用近场尺度 $5$ m；外包角域尖点引入的零首测距离假设不参与角度计算。同地点误差固定，故在生成 $Q$ 时剔除距已有检测站过近的候选，默认最小间距取 $25$ m。

\subsubsection{选点算法及第二次反馈处理}

具体步骤如下，所用信息均来自题设先验和已经取得的示向度。
\begin{enumerate}
\item 计算首测外包多边形 $P_1$ 及其最小包围圆 $(C^*,R^*)$，按式\eqref{eq:q2-nonempty}检验保证再接收域是否非空。
\item 构造 $K_0$ 的安全内多边形，候选中加入 $C^*$、至多 $16$ 个内多边形顶点，以及 $C^*\pm h(-\sin\varphi_1,\cos\varphi_1)$，其中 $h\in\{50,150,300,500\}$ m；去重并剔除距首站不足 $25$ m 的点。
\item 对每个候选重新检查到 $P_1$ 全部顶点的距离，保留满足式\eqref{eq:q2-safe}者；用至多 $12$ 个抽样源域顶点及最小圆心组成 $H$。
\item 计算 $J(q)$，选择有限集合 $Q$ 中评分最小的检测点；并列时依次选择移动距离较短、坐标字典序较小的点。
\item 在选定点测量。若返回示向度，则将第二次误差角域、接收距离上界与原区域求交，重新计算 $D$ 和 $R^*$；若返回近场响应，则在当前点进行光学定位。后验区域是否足够小以实际反馈为准，不由预测评分代替。
\end{enumerate}
该算法保证保留候选的再次接收，并在\textbf{有限候选集内}择优。顶点抽样、误差线性化和分母截断使它不构成连续检测域上的全局最优算法。若程序因数据不一致或离散筛选未获得合格候选，应报告这一状态，不能把任意位置标记为保证接收点。

\subsubsection{选点算例与定位效果}

取 $S_1=(0,0)$、首测 $\theta_1=0^\circ$。选点算法只使用首测信息；为比较第二次测量后的区域，另取真实源 $G=(1000,0)$ 构造一个受控算例。$G$ 仅用于生成观测及事后检验，没有传入选点函数。第二次示向采用无附加误差的真实方位并舍入至两位小数，区域计算采用 $1.005^\circ$ 保守半角。

\begin{table}[htbp]
\centering
\caption{首测条件相同的第二检测点对照}\label{tab:q2-example}
\small
\begin{tabular}{lccc}
\toprule
方案 & 第二检测点 / m & 第二示向 & 后验 $R^*$ / m \\
\midrule
仅有首测 & --- & --- & 750.2308 \\
共线前进 & $(500,0)$ & $0.00^\circ$ & 500.1539 \\
主动选点 & $(857.728610,514.102744)$ & $285.47^\circ$ & 22.8123 \\
\bottomrule
\end{tabular}
\end{table}

本例得到 $25$ 个合格候选，所选点通过全部顶点的 $1000$ m 距离核验。表\ref{tab:q2-example}表明，共线前进仍保留较大的纵向不确定性，而横向选点显著改善交会几何。两方案分别移动 $500$ m 和 $1000$ m，故此表用于说明几何缩域效果，不能据此宣称等成本最优或完整清除任务的时间优势。

\begin{figure}[htbp]
\centering
\begin{tikzpicture}[x=0.006cm,y=0.006cm,font=\small]
\draw[->,gray] (-100,0)--(1650,0) node[right] {$x$/m};
\draw[->,gray] (0,-740)--(0,800) node[above] {$y$/m};
\filldraw[fill=green!12,draw=green!45!black] (998.795456,-49.067674) -- (998.795456,49.067674) -- (989.176510,146.730474) -- (970.031253,242.980180) -- (941.544065,336.889853) -- (903.989293,427.555093) -- (857.728610,514.102744) -- (803.207531,595.699304) -- (760.969704,647.166247) -- (759.048875,645.245418) -- (696.792469,569.385767) -- (642.271390,487.789207) -- (596.010707,401.241556) -- (558.455935,310.576316) -- (529.968747,216.666643) -- (510.823490,120.416937) -- (501.204544,22.754137) -- (501.204544,-22.754137) -- (510.823490,-120.416937) -- (529.968747,-216.666643) -- (558.455935,-310.576316) -- (596.010707,-401.241556) -- (642.271390,-487.789207) -- (696.792469,-569.385767) -- (759.048875,-645.245418) -- (760.969704,-647.166247) -- (803.207531,-595.699304) -- (857.728610,-514.102744) -- (903.989293,-427.555093) -- (941.544065,-336.889853) -- (970.031253,-242.980180) -- (989.176510,-146.730474) -- cycle;
\filldraw[fill=blue!22,draw=blue!60!black,line width=.7pt] (1500.000000,26.313537) -- (0.000000,0.000000) -- (1500.000000,-26.313537) -- cycle;
\draw[densely dashed,->,red!70!black] (0,0)--(857.728610,514.102744);
\fill (0,0) circle[radius=1.7pt] node[below left] {$S_1$};
\fill[red!70!black] (857.728610,514.102744) circle[radius=1.8pt] node[above right] {$S_2$};
\fill (500,0) circle[radius=1.5pt] node[below] {共线对照点};
\fill[orange!80!black] (1000,0) circle[radius=1.6pt] node[above right] {$G$（仅评价）};
\node[green!30!black] at (750,-340) {$K_0$ 的内近似};
\node[blue!70!black] at (1370,170) {$P_1$};
\draw[blue!60!black,->] (1370,135)--(1370,23);
\foreach \x in {500,1000,1500}{\draw[gray] (\x,0)--(\x,-18);}
\end{tikzpicture}
\caption{首测源候选区域与第二检测点的安全候选域。坐标等比例，窄扇形按实际半角绘制；绿色区域为保证再接收域的内近似。}\label{fig:q2-domain}
\end{figure}

还应注意，主动选点后的 $R^*=22.8123$ m 仍大于 $20$ m，因此本例尚不能保证一次光学清除，需要继续观测或使用后续问题的有限覆盖清除策略。该算例也不意味着任意源位置和误差组合都能在两次测向后达到相同精度。问题二得到的是可执行的检测点候选域和选取方法，其有效性最终由实际第二次观测形成的定位区域检验。
